\section{Целочисленные предикаты: cross, dot, ориентация}
Вычисляет векторное и скалярное произведения и знак поворота для точек \texttt{pair<ll,ll>}; положительный поворот направлен против часовой стрелки, нулевой означает коллинеарность.
Исходные координаты до вычитания должны удовлетворять $|x|,|y|\le10^9$, иначе произведения могут переполнить \texttt{ll}.
Нужны стандартные заголовки, \texttt{std} и \texttt{ll = long long}; время и память $O(1)$.
\begin{lstlisting}[name=geometry_integer]
ll cross_integer(pair<ll, ll> a, pair<ll, ll> b) {
    return a.first * b.second - a.second * b.first;
}

ll dot_integer(pair<ll, ll> a, pair<ll, ll> b) {
    return a.first * b.first + a.second * b.second;
}

int orientation(pair<ll, ll> a, pair<ll, ll> b, pair<ll, ll> c) {
    ll value = cross_integer({b.first - a.first, b.second - a.second},
                             {c.first - a.first, c.second - a.second});
    return (value > 0) - (value < 0);
}
\end{lstlisting}
% Точные знаки поворота и скалярные произведения за $O(1)$.
% Точки и векторы --- \texttt{pair<ll,ll>}, положительный поворот --- против
% часовой стрелки. Во всех листингах нужны стандартные заголовки,
% \texttt{using namespace std}, внешние \texttt{ll = long long} и
% \texttt{db = double}. Глобальных типов и констант здесь нет.
% Для исходных целочисленных координат $|x|,|y|\le10^9$ разности не больше
% $2\cdot10^9$, каждое произведение не больше $4\cdot10^{18}$, сумма или
% разность двух произведений не больше $8\cdot10^{18}<2^{63}-1$.
% Поэтому даже cross/dot двух разностей безопасны в \texttt{ll}; расширенная
% целочисленная арифметика не требуется. Для больших координат гарантии нет.
%
% Нулевой знак включает совпадающие точки. Ограничение относится к координатам
% до вычитания; суммы Минковского могут уже выйти за этот диапазон.

\section{Точка на отрезке и пересечение отрезков}
Проверяет принадлежность точки замкнутому отрезку и пересечение двух замкнутых отрезков, возвращая логический ответ.
Касания концами и коллинеарное перекрытие считаются пересечением, вырожденный отрезок допустим.
Нужна \texttt{orientation} с исходными целочисленными координатами по модулю не больше $10^9$; время и память $O(1)$.
\begin{lstlisting}[name=geometry_segments]
bool on_segment(pair<ll, ll> p, pair<ll, ll> a, pair<ll, ll> b) {
    return orientation(a, b, p) == 0 &&
        min(a.first, b.first) <= p.first && p.first <= max(a.first, b.first) &&
        min(a.second, b.second) <= p.second && p.second <= max(a.second, b.second);
}

bool segments_intersect(pair<ll, ll> a, pair<ll, ll> b,
                        pair<ll, ll> c, pair<ll, ll> d) {
    int ab_c = orientation(a, b, c), ab_d = orientation(a, b, d);
    int cd_a = orientation(c, d, a), cd_b = orientation(c, d, b);
    if (ab_c == 0 && on_segment(c, a, b)) return true;
    if (ab_d == 0 && on_segment(d, a, b)) return true;
    if (cd_a == 0 && on_segment(a, c, d)) return true;
    if (cd_b == 0 && on_segment(b, c, d)) return true;
    return ab_c * ab_d < 0 && cd_a * cd_b < 0;
}
\end{lstlisting}
% Точные замкнутые предикаты за $O(1)$: касание концами и коллинеарное
% перекрытие считаются пересечением. Отрезок может быть точкой.
% Использует \texttt{orientation} и его ограничения координат.
%
% Предикат не восстанавливает точку или общий отрезок. Для координат
% пересечения непараллельных прямых ниже есть вещественная функция.

\section{Вещественные прямые: пересечение, проекция, отражение}
Для точек \texttt{pair<db,db>} находит пересечение прямых, проекцию и отражение точки; прямая задаётся двумя различными точками.
Пересечение возвращает \texttt{nullopt} для параллельных или совпадающих прямых, проекция и отражение --- для вырожденной прямой.
Здесь и далее \texttt{db = double}, координаты и промежуточные вычисления должны быть конечны, а \texttt{eps} --- конечным и неотрицательным; допуск подбирается под задачу и не гарантирует точности на почти вырожденных данных.
Время и память $O(1)$; параллельность проверяется относительным угловым допуском, совпадение концов --- точно.
\begin{lstlisting}[name=geometry_lines]
db cross_real(pair<db, db> a, pair<db, db> b, pair<db, db> c) {
    return (b.first - a.first) * (c.second - a.second) -
           (b.second - a.second) * (c.first - a.first);
}

db squared_distance(pair<db, db> a, pair<db, db> b) {
    db x = a.first - b.first, y = a.second - b.second;
    return x * x + y * y;
}

optional<pair<db, db>> line_intersection(
        pair<db, db> a, pair<db, db> b,
        pair<db, db> c, pair<db, db> d, db eps) {
    db ux = b.first - a.first, uy = b.second - a.second;
    db vx = d.first - c.first, vy = d.second - c.second;
    db determinant = ux * vy - uy * vx;
    if (a == b || c == d ||
        abs(determinant) <= eps * hypot(ux, uy) * hypot(vx, vy)) return nullopt;
    db t = ((c.first - a.first) * vy - (c.second - a.second) * vx) / determinant;
    return pair<db, db>{a.first + t * ux, a.second + t * uy};
}

optional<pair<db, db>> project_line(
        pair<db, db> p, pair<db, db> a, pair<db, db> b) {
    if (a == b) return nullopt;
    db x = b.first - a.first, y = b.second - a.second;
    db t = ((p.first - a.first) * x + (p.second - a.second) * y) / (x*x + y*y);
    return pair<db, db>{a.first + t*x, a.second + t*y};
}

optional<pair<db, db>> reflect_line(
        pair<db, db> p, pair<db, db> a, pair<db, db> b) {
    auto projection = project_line(p, a, b);
    if (!projection) return nullopt;
    return pair<db, db>{2*projection->first - p.first, 2*projection->second - p.second};
}
\end{lstlisting}
% Операции за $O(1)$ с точками \texttt{pair<db,db>}. Прямая задаётся двумя
% различными точками. Пересечение возвращает \texttt{nullopt} для параллельных
% и совпадающих прямых; проекция и отражение --- для вырожденной прямой.
% Здесь и далее \texttt{eps} --- конечное неотрицательное число. Все входные
% вещественные координаты конечны; промежуточные произведения, квадраты и
% деления должны оставаться конечными. При \texttt{double} практически
% разумны координаты порядка $10^6$ и \texttt{eps} порядка $10^{-9}$,
% но это не гарантия точных предикатов: почти параллельные прямые,
% почти касания и плохо обусловленные данные требуют выбора допуска под задачу.
% \texttt{cross\_real(a,b,c)} считает $(b-a)\times(c-a)$.
% Параллельность проверяется относительным угловым допуском; совпадение
% концов прямой --- точным равенством. Сортировки всегда без допуска.
%
% \texttt{nullopt} намеренно не различает параллельность и совпадение.
% Для точных рациональных пересечений нужен другой числовой контракт.
% Общие вещественные ограничения этого блока действуют во всех следующих.

\section{Расстояния от точки до прямой и отрезка}
Возвращает евклидово расстояние от точки до прямой или замкнутого отрезка, заданных двумя точками.
Для прямой из одной точки возвращает \texttt{nullopt}, для отрезка из одной точки --- расстояние до неё.
Нужны \texttt{cross\_real} и \texttt{squared\_distance} с конечной вещественной арифметикой; время и память $O(1)$.
\begin{lstlisting}[name=geometry_distances]
optional<db> distance_line(pair<db, db> p, pair<db, db> a, pair<db, db> b) {
    if (a == b) return nullopt;
    return abs(cross_real(a, b, p)) / hypot(b.first-a.first, b.second-a.second);
}

db distance_segment(pair<db, db> p, pair<db, db> a, pair<db, db> b) {
    if (a == b) return sqrt(squared_distance(p, a));
    db x = b.first-a.first, y = b.second-a.second;
    db t = clamp(((p.first-a.first)*x + (p.second-a.second)*y) / (x*x+y*y), 0.0, 1.0);
    return hypot(p.first-a.first-t*x, p.second-a.second-t*y);
}
\end{lstlisting}
% Евклидовы расстояния за $O(1)$. Для прямой из одной точки возвращается
% \texttt{nullopt}, для отрезка из одной точки --- расстояние до неё.
% Использует \texttt{cross\_real} и \texttt{squared\_distance}.
%
% Для сравнения расстояний можно работать с квадратами без извлечения корня.
% Почти нулевой, но ненулевой отрезок здесь не заменяется точкой автоматически.

\section{Ориентированная площадь многоугольника}
По циклическому списку вершин возвращает ориентированную площадь: положительную при обходе против часовой стрелки, для обычной площади возьмите модуль.
Пустой список, точка и отрезок дают ноль; для самопересечений результат алгебраический, а не площадь объединения.
Нужна \texttt{cross\_real} с конечной вещественной арифметикой; время $O(n)$, дополнительная память $O(1)$.
\begin{lstlisting}[name=geometry_area]
db polygon_area(const vector<pair<db, db>>& polygon) {
    db twice = 0;
    for (int i = 1; i + 1 < (int)polygon.size(); ++i)
        twice += cross_real(polygon[0], polygon[i], polygon[i+1]);
    return twice / 2;
}
\end{lstlisting}
% Площадь за $O(n)$ и $O(1)$ памяти; знак положителен для обхода против часовой
% стрелки. Пустой список, точка и отрезок имеют нулевую площадь.
% Использует \texttt{cross\_real}. Суммирование треугольников относительно
% первой вершины уменьшает потерю точности при большом общем смещении.
%
% Для обычной площади берётся модуль; для самопересечений результат
% алгебраический, не площадь объединения. Даже при целых координатах
% сумма по всем рёбрам не обязана помещаться в \texttt{ll}.

\section{Точка в произвольном простом многоугольнике}
По целочисленным вершинам простого многоугольника и точке возвращает $-1$ снаружи, $0$ на границе, $1$ внутри за $O(n)$.
Направление обхода произвольно, первую вершину в конце не повторяют; пустой список не содержит точек, одна или две точки задают только границу.
Нужны \texttt{orientation} и \texttt{on\_segment}, исходные координаты по модулю не больше $10^9$.
\begin{lstlisting}[name=geometry_point_polygon]
int point_polygon(const vector<pair<ll, ll>>& polygon, pair<ll, ll> p) {
    bool inside = false;
    int n = polygon.size();
    for (int i = 0; i < n; ++i) {
        auto a = polygon[i], b = polygon[(i+1)%n];
        if (on_segment(p, a, b)) return 0;
        if ((a.second > p.second) != (b.second > p.second)) {
            int turn = orientation(a, b, p);
            if ((turn > 0) == (b.second > a.second)) inside = !inside;
        }
    }
    return inside ? 1 : -1;
}
\end{lstlisting}
% За $O(n)$ возвращает $-1$ снаружи, $0$ на границе, $1$ внутри.
% Целочисленные координаты, направление обхода произвольно. Пустой список
% не содержит точек; списки из одной или двух точек трактуются как граница.
% Использует \texttt{orientation} и \texttt{on\_segment} с их ограничениями.
% Луч направлен вправо, полуоткрытая проверка высот не удваивает вершины.
%
% Последнюю вершину не нужно повторять. При самопересечениях получается
% правило чётности, а не проверка принадлежности простому многоугольнику.

\section{Точка в выпуклом многоугольнике за логарифм}
Проверяет точку в выпуклом многоугольнике за $O(\log n)$ и возвращает $-1$ снаружи, $0$ на границе, $1$ внутри.
При $n\ge3$ вершины должны идти против часовой стрелки без повторов и промежуточных коллинеарных точек; начальная вершина любая, пустой список, точка и отрезок допустимы.
Нужны целочисленные предикаты с исходными координатами по модулю не больше $10^9$; дополнительная память $O(1)$, выпуклость входа не проверяется.
\begin{lstlisting}[name=geometry_point_convex]
int point_convex(const vector<pair<ll, ll>>& polygon, pair<ll, ll> p) {
    int n = polygon.size();
    if (n == 0) return -1;
    if (n == 1) return p == polygon[0] ? 0 : -1;
    if (n == 2) return on_segment(p, polygon[0], polygon[1]) ? 0 : -1;
    int first = orientation(polygon[0], polygon[1], p);
    int last = orientation(polygon[0], polygon[n-1], p);
    if (first < 0 || last > 0) return -1;
    if (first == 0) return on_segment(p, polygon[0], polygon[1]) ? 0 : -1;
    if (last == 0) return on_segment(p, polygon[0], polygon[n-1]) ? 0 : -1;
    int left = 1, right = n-1;
    while (right-left > 1) {
        int middle = (left+right)/2;
        if (orientation(polygon[0], polygon[middle], p) >= 0) left = middle;
        else right = middle;
    }
    return orientation(polygon[left], polygon[right], p);
}
\end{lstlisting}
% Та же классификация $-1/0/1$ за $O(\log n)$ и $O(1)$ памяти.
% При $n\ge3$ нужен строго выпуклый многоугольник против часовой стрелки,
% без повторов и промежуточных коллинеарных вершин. Начальная вершина любая.
% Использует целочисленные предикаты. Пустой список, точка и отрезок допустимы.
%
% Очистку и разворот входа следует делать один раз до запросов, а не внутри
% функции: иначе логарифмическое время потеряется. Проверки контракта нет.

\iffalse
\section{Диаметр выпуклого многоугольника: вращающиеся калиперы}
По циклическому списку вершин выпуклого многоугольника возвращает квадрат максимального расстояния за $O(n)$ времени и $O(1)$ дополнительной памяти.
При $n\ge3$ многоугольник должен быть строго выпуклым и обходиться против часовой стрелки; пустой список и одна точка дают $0$, две точки --- квадрат расстояния между ними.
Нужны \texttt{cross\_real} и \texttt{squared\_distance} с конечной вещественной арифметикой; почти коллинеарные данные могут нарушить корректность из-за округления.
\begin{lstlisting}[name=geometry_diameter]
db convex_diameter_squared(const vector<pair<db, db>>& polygon) {
    int n = polygon.size();
    if (n < 2) return 0;
    if (n == 2) return squared_distance(polygon[0], polygon[1]);
    db answer = 0;
    int opposite = 1;
    for (int i = 0; i < n; ++i) {
        int next = (i+1)%n;
        while (cross_real(polygon[i], polygon[next], polygon[(opposite+1)%n]) >
               cross_real(polygon[i], polygon[next], polygon[opposite]))
            opposite = (opposite+1)%n;
        answer = max({answer, squared_distance(polygon[i], polygon[opposite]),
                      squared_distance(polygon[next], polygon[opposite])});
    }
    return answer;
}
\end{lstlisting}
\fi

% Квадрат максимального расстояния между вершинами за $O(n)$ и $O(1)$ памяти.
% Для пустого списка и одной точки результат $0$, для двух точек --- их
% квадрат расстояния. При $n\ge3$ вход строго выпуклый, против часовой стрелки.
% Использует \texttt{cross\_real} и \texttt{squared\_distance}. Равные площади
% не продвигают указатель; обе вершины текущего ребра проверяются.
%
% Можно сохранять пару индексов при обновлении ответа. Численная ошибка
% на почти коллинеарном входе может нарушить предпосылки калиперов.

\iffalse
\section{Отсечение многоугольника полуплоскостью}
Отсекает выпуклый многоугольник левой замкнутой полуплоскостью ориентированной прямой $a\to b$ и возвращает циклический список вершин, возможно пустой, точку или отрезок.
Нужны $a\ne b$, \texttt{cross\_real} и конечная вещественная арифметика; время и память $O(n)$.
Допуск \texttt{eps} измеряется в единицах расстояния: граница сдвигается до $(b-a)\times(p-a)\ge-\texttt{eps}\,|b-a|$, и пересечения вычисляются с этой же границей.
Для невыпуклого входа один список не позволяет представить несвязное пересечение.
\begin{lstlisting}[name=geometry_clipping]
vector<pair<db, db>> clip_halfplane(const vector<pair<db, db>>& polygon,
        pair<db, db> a, pair<db, db> b, db eps) {
    assert(a != b && eps >= 0);
    vector<pair<db, db>> result;
    db shift = eps * hypot(b.first-a.first, b.second-a.second);
    auto append = [&](pair<db, db> p) {
        if (result.empty() || result.back() != p) result.push_back(p);
    };
    int n = polygon.size();
    for (int i = 0; i < n; ++i) {
        auto p = polygon[i], q = polygon[(i+1)%n];
        db from = cross_real(a, b, p) + shift;
        db to = cross_real(a, b, q) + shift;
        if (from >= 0) append(p);
        if ((from >= 0) != (to >= 0)) {
            db t = from / (from-to);
            append({p.first + t*(q.first-p.first), p.second + t*(q.second-p.second)});
        }
    }
    if (result.size() > 1 && result.front() == result.back()) result.pop_back();
    return result;
}
\end{lstlisting}
\fi

% Пересечение выпуклого многоугольника с левой замкнутой полуплоскостью
% ориентированной прямой $a\to b$ за $O(n)$ времени и памяти.
% Результат --- циклический список, возможно пустой, точка или отрезок.
% Для невыпуклого входа один список не представляет несвязный результат.
% Использует \texttt{cross\_real}. Требуется $a\ne b$. Допуск имеет единицы
% расстояния: реально отсекается по сдвинутой границе
% $(b-a)\times(p-a)\ge-\texttt{eps}\,|b-a|$.
% Пересечения считаются с этой же границей, а не с исходной прямой.
% Удаляются точные соседние повторы, близкие точки не склеиваются.
%
% \texttt{eps=0} не сдвигает геометрию, но не устраняет ошибки округления.
% При положительном допуске площадь и вершины относятся к расширенной области.

\iffalse
\section{Ограниченное пересечение полуплоскостей отсечениями}
\subsection*{Для чего нужен}
Честный простой вариант: пересечение $k$ левых полуплоскостей с явно
заданным квадратом $[-B,B]^2$. Это не алгоритм за $O(k\log k)$:
время $O(k(k+4))$, память $O(k+4)$, так как каждый шаг добавляет не более
одной вершины выпуклого многоугольника в точной арифметике.
\subsection*{Реализация}
Использует \texttt{clip\_halfplane}. $B>0$ конечен; концы каждой прямой
различны. \texttt{eps} передаётся отсечению и расширяет полуплоскости,
но не bounding box. Пустой набор ограничений возвращает квадрат.
\begin{lstlisting}[name=geometry_halfplanes]
vector<pair<db, db>> bounded_halfplanes(
        const vector<pair<pair<db, db>, pair<db, db>>>& halfplanes,
        db bound, db eps) {
    assert(isfinite(bound) && bound > 0 && eps >= 0);
    vector<pair<db, db>> polygon = {
        {-bound, -bound}, {bound, -bound}, {bound, bound}, {-bound, bound}};
    for (const auto& [a, b] : halfplanes) {
        assert(a != b);
        polygon = clip_halfplane(polygon, a, b, eps);
    }
    return polygon;
}
\end{lstlisting}
\subsection*{Возможные улучшения и дополнения}
Функция не определяет неограниченность исходного пересечения. Bounding box
является частью задачи; чтобы получить исходную ограниченную область,
нужно заранее доказать, что она целиком лежит внутри квадрата.
Плохая обусловленность может нарушить теоретическую структуру результата.


\fi
\iffalse
\section{Сумма Минковского выпуклых многоугольников}
\subsection*{Для чего нужен}
Граница множества $A+B$ за $O(n+m)$ времени и памяти слиянием направлений
рёбер. Непустые входы при размере от трёх строго выпуклы, идут против часовой
стрелки и не повторяют начало. Точки и отрезки из разных концов допустимы.
\subsection*{Реализация}
Использует \texttt{cross\_real}. Пустой аргумент даёт пустую сумму.
Начало переносится в минимальную по $(y,x)$ вершину. Равные направления
объединяются по точному нулю cross; численно почти параллельные рёбра
могут оставить почти коллинеарные вершины. Результат начинается снизу слева.
\begin{lstlisting}[name=geometry_minkowski]
vector<pair<db, db>> minkowski_sum(vector<pair<db, db>> a,
                                   vector<pair<db, db>> b) {
    if (a.empty() || b.empty()) return {};
    auto rotate_lowest = [](vector<pair<db, db>>& polygon) {
        auto lowest = min_element(polygon.begin(), polygon.end(), [](auto p, auto q) {
            return pair{p.second, p.first} < pair{q.second, q.first};
        });
        rotate(polygon.begin(), lowest, polygon.end());
    };
    rotate_lowest(a);
    rotate_lowest(b);
    if (a.size() == 1 || b.size() == 1) {
        if (b.size() != 1) swap(a, b);
        for (auto& p : a) { p.first += b[0].first; p.second += b[0].second; }
        return a;
    }
    int n = a.size(), m = b.size(), i = 0, j = 0;
    vector<pair<db, db>> result;
    while (i < n || j < m) {
        result.push_back({a[i%n].first+b[j%m].first, a[i%n].second+b[j%m].second});
        if (i == n) { ++j; continue; }
        if (j == m) { ++i; continue; }
        pair<db, db> u = {a[(i+1)%n].first-a[i%n].first,
                         a[(i+1)%n].second-a[i%n].second};
        pair<db, db> v = {b[(j+1)%m].first-b[j%m].first,
                         b[(j+1)%m].second-b[j%m].second};
        db turn = cross_real({0, 0}, u, v);
        if (turn > 0) ++i;
        else if (turn < 0) ++j;
        else { ++i; ++j; }
    }
    return result;
}
\end{lstlisting}
\subsection*{Возможные улучшения и дополнения}
Нельзя произвольно подавать невыпуклый вход или обход по часовой стрелке.
Для исходных произвольных наборов сначала нужна выпуклая оболочка;
её построение не входит в линейную оценку этой функции.


\fi
\iffalse
\section{Ближайшая пара точек}
\subsection*{Для чего нужен}
Минимальный квадрат расстояния за $O(n\log n)$ времени и $O(n)$ памяти.
Для менее двух точек возвращается \texttt{nullopt}, повторы дают $0$.
Вход копируется; исходный порядок не изменяется.
\subsection*{Реализация}
Использует \texttt{squared\_distance}. После рекурсии диапазон отсортирован
по $y$; слияние и просмотр полосы линейны благодаря упаковке точек.
Все квадраты расстояний должны помещаться в конечный \texttt{db}.
\begin{lstlisting}[name=geometry_closest]
optional<db> closest_pair_squared(vector<pair<db, db>> points) {
    int n = points.size();
    if (n < 2) return nullopt;
    sort(points.begin(), points.end());
    if (adjacent_find(points.begin(), points.end()) != points.end()) return 0;
    vector<pair<db, db>> buffer(n);
    auto by_y = [](auto a, auto b) {
        return pair{a.second, a.first} < pair{b.second, b.first};
    };
    function<db(int, int)> solve = [&](int left, int right) -> db {
        db answer = numeric_limits<db>::infinity();
        if (right-left <= 3) {
            for (int i = left; i < right; ++i)
                for (int j = i+1; j < right; ++j)
                    answer = min(answer, squared_distance(points[i], points[j]));
            sort(points.begin()+left, points.begin()+right, by_y);
            return answer;
        }
        int middle = (left+right)/2;
        db split = points[middle].first;
        db first = solve(left, middle), second = solve(middle, right);
        answer = min(first, second);
        merge(points.begin()+left, points.begin()+middle,
              points.begin()+middle, points.begin()+right, buffer.begin()+left, by_y);
        copy(buffer.begin()+left, buffer.begin()+right, points.begin()+left);
        int count = 0;
        for (int i = left; i < right; ++i) {
            db dx = points[i].first-split;
            if (dx*dx >= answer) continue;
            for (int j = count-1; j >= 0; --j) {
                db dy = points[i].second-buffer[j].second;
                if (dy*dy >= answer) break;
                answer = min(answer, squared_distance(points[i], buffer[j]));
            }
            buffer[count++] = points[i];
        }
        return answer;
    };
    return solve(0, n);
}
\end{lstlisting}
\subsection*{Возможные улучшения и дополнения}
Можно хранить исходные индексы вместе с точками для восстановления пары.
Здесь возвращается только расстояние, без добавления отдельного типа точки.


\fi
\iffalse
\section{Пересечение окружности и прямой}
По центру и радиусу окружности и двум точкам прямой возвращает от нуля до двух точек пересечения за $O(1)$.
Нужны конечный $r\ge0$, различные точки прямой, \texttt{project\_line} и \texttt{squared\_distance}; нулевой радиус допустим.
При $|d-r|\le\texttt{eps}$ возвращается одна точка --- проекция центра на прямую, то есть приближённое касание даже без точного пересечения.
\begin{lstlisting}[name=geometry_circle_line]
vector<pair<db, db>> circle_line(pair<db, db> center, db radius,
        pair<db, db> a, pair<db, db> b, db eps) {
    assert(radius >= 0 && a != b && eps >= 0);
    auto foot = *project_line(center, a, b);
    db distance = sqrt(squared_distance(center, foot));
    if (distance > radius+eps) return {};
    if (abs(distance-radius) <= eps) return {foot};
    db length = hypot(b.first-a.first, b.second-a.second);
    db offset = sqrt(max(0.0, (radius-distance)*(radius+distance)));
    db x = (b.first-a.first)/length*offset;
    db y = (b.second-a.second)/length*offset;
    return {{foot.first-x, foot.second-y}, {foot.first+x, foot.second+y}};
}
\end{lstlisting}
\fi

% От нуля до двух точек за $O(1)$. Радиус $r\ge0$ конечен, прямая $a\ne b$.
% Нулевой радиус допустим; касание возвращает ровно одну точку.
% Использует \texttt{project\_line} и \texttt{squared\_distance}.
% Допуск абсолютный по расстоянию: при $|d-r|\le\texttt{eps}$ возвращается
% проекция центра (приближённое касание), даже если точного касания нет.
%
% Для отрезка найденные точки дополнительно фильтруются по параметру.
% Порядок двух точек следует направлению прямой.

\iffalse
\section{Пересечение двух окружностей}
Находит общие точки двух окружностей за $O(1)$: пустой вектор означает отсутствие пересечений, \texttt{nullopt} --- бесконечно много точек.
Радиусы должны быть конечными и неотрицательными; одинаковые окружности нулевого радиуса имеют ровно одну общую точку.
Нужна \texttt{squared\_distance}; совпадение центров и радиусов проверяется точно, а \texttt{eps} задаёт допуск по расстоянию для разделения и касания, поэтому точки касания могут быть приближёнными.
\begin{lstlisting}[name=geometry_circle_circle]
optional<vector<pair<db, db>>> circle_circle(pair<db, db> a, db r,
        pair<db, db> b, db s, db eps) {
    assert(r >= 0 && s >= 0 && eps >= 0);
    vector<pair<db, db>> result;
    db distance = sqrt(squared_distance(a, b));
    if (a == b) {
        if (r != s) return result;
        if (r == 0) return vector<pair<db, db>>{a};
        return nullopt;
    }
    if (distance > r+s+eps || distance < abs(r-s)-eps) return result;
    db along = (distance + (r-s)*(r+s)/distance)/2;
    pair<db, db> unit = {(b.first-a.first)/distance, (b.second-a.second)/distance};
    pair<db, db> base = {a.first+along*unit.first, a.second+along*unit.second};
    if (abs(distance-r-s) <= eps || abs(distance-abs(r-s)) <= eps)
        return vector<pair<db, db>>{base};
    db height = sqrt(max(0.0, (r-along)*(r+along)));
    result.push_back({base.first-height*unit.second, base.second+height*unit.first});
    result.push_back({base.first+height*unit.second, base.second-height*unit.first});
    return result;
}
\end{lstlisting}
\fi

% За $O(1)$ возвращает \texttt{optional<vector<pair<db,db>>>}:
% \texttt{nullopt} означает бесконечно много точек, пустой вектор --- ни одной.
% Радиусы конечны и неотрицательны; одинаковые окружности нулевого радиуса
% имеют ровно одну общую точку, а не бесконечность.
% Использует \texttt{squared\_distance}. Совпадение центров и радиусов
% проверяется точно; \texttt{eps} применяется к расстояниям при разделении
% и касании, но не превращает близкие окружности в совпадающие.
%
% При приближённом касании возвращаемая точка лишь приближённо лежит на обеих
% окружностях. Очень близкие центры делают вычисление основания неустойчивым.

\iffalse
\section{Общие касательные двух окружностей}
\subsection*{Для чего нужен}
От нуля до четырёх прямых за $O(1)$. Радиусы строго положительны.
Прямая представлена парой \texttt{(точка, единичное направление)}, а не
двумя точками касания: при касании окружностей точки касания совпадают.
\texttt{nullopt} --- бесконечное семейство для совпадающих окружностей.
\subsection*{Реализация}
Использует \texttt{squared\_distance}. Совпадение окружностей точное.
Для каждого знака радиуса решается уравнение единичной нормали;
\texttt{eps} --- абсолютный допуск по расстоянию между центрами.
Вблизи касания нормаль нормируется, поэтому отклонение от второго радиуса
не превосходит допуска с учётом округления.
\begin{lstlisting}[name=geometry_tangents]
optional<vector<pair<pair<db, db>, pair<db, db>>>> common_tangents(
        pair<db, db> a, db r, pair<db, db> b, db s, db eps) {
    assert(r > 0 && s > 0 && eps >= 0);
    vector<pair<pair<db, db>, pair<db, db>>> result;
    if (a == b) {
        if (r == s) return nullopt;
        return result;
    }
    db distance = sqrt(squared_distance(a, b));
    db x = (b.first-a.first)/distance, y = (b.second-a.second)/distance;
    for (int sign : {1, -1}) {
        db difference = r-sign*s;
        if (distance < abs(difference)-eps) continue;
        db cosine = clamp(difference/distance, -1.0, 1.0);
        db sine = sqrt(max(0.0, (1-cosine)*(1+cosine)));
        if (abs(distance-abs(difference)) <= eps) {
            cosine = difference >= 0 ? 1 : -1;
            sine = 0;
        }
        for (int side : {-1, 1}) {
            db nx = x*cosine-side*y*sine, ny = y*cosine+side*x*sine;
            result.push_back({{a.first+r*nx, a.second+r*ny}, {-ny, nx}});
            if (sine == 0) break;
        }
    }
    return result;
}
\end{lstlisting}
\subsection*{Возможные улучшения и дополнения}
Нулевые радиусы здесь намеренно исключены: касательные к точкам требуют
отдельного соглашения о совпадающих ответах. Для проверки прямой достаточно
расстояния от каждого центра до неё.


\fi
\iffalse
\section{Минимальная охватывающая окружность}
\subsection*{Для чего нужен}
Рандомизированный инкрементальный алгоритм: ожидаемое время $O(n)$,
худшее $O(n^3)$, память $O(n)$ из-за копии входа. Результат ---
\texttt{(центр, радиус)}, пустой вход даёт \texttt{nullopt}.
Одна точка и повторы допустимы, коллинеарные наборы тоже.
\subsection*{Реализация}
Использует \texttt{squared\_distance}; генератор передаётся снаружи.
Допуск абсолютный по радиусу: точка внутри, если расстояние не больше
$r+\texttt{eps}$. Поэтому при положительном допуске результат приближённый.
Определитель и центр окружности трёх точек считаются в \texttt{long double};
точно нулевой определитель заменяется окружностью на самой дальней паре.
Почти коллинеарные тройки остаются плохо обусловленными; результат и
промежуточная арифметика обязаны быть конечными.
\begin{lstlisting}[name=geometry_enclosing]
optional<pair<pair<db, db>, db>> minimum_enclosing_circle(
        vector<pair<db, db>> points, mt19937_64& random, db eps) {
    assert(eps >= 0);
    if (points.empty()) return nullopt;
    shuffle(points.begin(), points.end(), random);
    pair<db, db> center = points[0];
    db radius = 0;
    auto outside = [&](pair<db, db> p) {
        return hypot(p.first-center.first, p.second-center.second) > radius+eps;
    };
    auto diameter = [&](pair<db, db> a, pair<db, db> b) {
        center = {a.first+(b.first-a.first)/2, a.second+(b.second-a.second)/2};
        radius = sqrt(squared_distance(a, b))/2;
    };
    for (int i = 0; i < (int)points.size(); ++i) {
        if (!outside(points[i])) continue;
        center = points[i];
        radius = 0;
        for (int j = 0; j < i; ++j) {
            if (!outside(points[j])) continue;
            diameter(points[i], points[j]);
            for (int k = 0; k < j; ++k) {
                if (!outside(points[k])) continue;
                auto a = points[i], b = points[j], c = points[k];
                long double bx = (long double)b.first-a.first;
                long double by = (long double)b.second-a.second;
                long double cx = (long double)c.first-a.first;
                long double cy = (long double)c.second-a.second;
                long double determinant = 2*(bx*cy-by*cx);
                if (determinant == 0) {
                    auto u = a, v = b;
                    if (squared_distance(a, c) > squared_distance(u, v)) { u = a; v = c; }
                    if (squared_distance(b, c) > squared_distance(u, v)) { u = b; v = c; }
                    diameter(u, v);
                } else {
                    long double first = bx*bx+by*by, second = cx*cx+cy*cy;
                    center = {a.first+(db)((first*cy-second*by)/determinant),
                              a.second+(db)((bx*second-cx*first)/determinant)};
                    radius = hypot(center.first-a.first, center.second-a.second);
                }
            }
        }
    }
    return pair{center, radius};
}
\end{lstlisting}
\subsection*{Возможные улучшения и дополнения}
Фиксированное зерно удобно для воспроизводимости, но ожидаемая линейная
оценка предполагает случайную перестановку, независимую от входа.
Допуск не является доказательством точного включения всех точек.

\fi
